\documentclass[a4paper]{panl}
\usepackage{cite}
\usepackage{wrapfig}
\usepackage[dvips]{graphicx}
\usepackage{amssymb}
\usepackage{amsfonts}
\usepackage{amsmath}
\usepackage{longtable}
\usepackage{rotating}
\usepackage{lscape}
\usepackage{epsfig}
\usepackage{multirow}
\usepackage[usenames]{color}
\usepackage{colortbl}
\usepackage[english,russian]{babel}

%\originalTeX
\russianTeX
\begin{document}
% Journal sections (see http://pkp.jinr.ru/index.php/PEPAN_LETTERS/about/editorialPolicies#focusAndScope)
%\issuearea{Physics of Elementary Particles and Atomic Nuclei. Theory}
% or in Russian
%\issuearea{ФИЗИКА ЭЛЕМЕНТАРНЫХ ЧАСТИЦ И АТОМНОГО ЯДРА. ТЕОРИЯ}

\title{Solution of the kinetic equation taking into account resonant the Compton process in a magnetized medium \\ Решение кинетического уравнения с учетом резонанса в комптоновском процессе в замагниченной среде}
\maketitle
\authors{A.\,Yarkov $^{a}$\footnote{E-mail: a12l@mail.ru},
D.\,Rumyantsev$^{a}$\footnote{E-mail: rda@uniyar.ac.ru}}
\setcounter{footnote}{0}
\authors{А.А.\,Ярков$^{a,1}$,
Д.А.\,Румянцев$^{a,2}$}
\from{$^{a}$\, Yaroslavl Demidov State University}
\from{$^{a}$\, Ярославский государственный университет им. П.Г. Демидова}

\begin{abstract}
% Russian translation of the abstract
Рассмотрено решение кинетического уравнения для нахождения функции распределения фотонов двух возможных поляризаций в равновесной нерелятивистской $e^+e^-$ плазме в относительно сильном магнитном поле с учётом резонанса на виртуальном электроне.
Используя преобразование Лапласа и разложение функции распределения по полиномам Лежандра задача сведена к системе дифференциальных уравнений на коэффициенты этого разложения. Решение представлено в виде интеграла в комплексной плоскости. \
\vspace{0.2cm}

The solution of the kinetic equation for finding the distribution function of photons of two possible polarizations in an equilibrium nonrelativistic $e^+e^-$ plasma in a relatively strong magnetic field is considered taking into account resonance on a virtual electron.
Using the Laplace transform and the expansion of the distribution function in Legendre polynomials, the problem is reduced to a system of differential equations for the coefficients of this expansion. The solution is presented as an integral in the complex plane.
\end{abstract}
\vspace*{6pt}

\noindent
 PACS: 13.60.Fz; 52.25.Dg; 52.27.Ep;

\label{sec:intro}
\section*{Введение}
В настоящее время является установленным
фактом, что наличие магнитного поля в широком
классе астрофизических объектов представляет
типичную ситуацию для наблюдаемой Вселенной.
При этом масштаб индукции магнитного поля
может варьироваться в очень широких пределах:
от крупномасштабных ($\sim 100$ кпк) межгалактических магнитных полей $\sim 10^{-21}$ Гс \cite{D.Ryu:2012}, до полей, 
реализующихся в сценарии ротационного взрыва сверхновой. При этом особый 
интерес представляют объекты с полями масштаба так
называемого критического значения $B_e=m^2/e\approx4.41\cdot10^{13}$ Гс.
\footnote{ В работе используются естесственная система единиц, где $c=h=k_B=1$, m -- масса электрона,
	$e>0$ -- элементарный заряд} К ним, в частности, относятся
изолированные нейтронные звезды, включающие в
себя радиопульсары и так называемые магнитары,
обладающими магнитными полями с индукцией от $B\sim10^{12}$ Гс (радиопульсары) до $B\sim4\cdot10^{14}$ Гс (магнитары). 

Анализ спектров излучения радиопульсаров
и магнитаров свидетельствует также о наличии
электрон-позитронной плазмы в их магнитосферах
с концентрацией порядка значения концентрации
Голдрайха–Джулиана \cite{Goldreich:1969}:
\begin{equation}
n_{GJ}\approx 3 \cdot 10^{13} \text{см}^{-3} \frac{B}{100B_e} \frac{10 c}{P}\, ,
\end{equation}
где P -- период вращения нейтронной звезды.

Естественно ожидать, что такие экстремальные
условия будут оказывать существенное влияние на
квантовые процессы, где в конечном или начальном
состоянии могут присутствовать как электрически
заряженные, так и электрически нейтральные частицы, например, электроны и фотоны.

Одним из таких процессов является активно обсуждаемое в настоящее время в литературе комптоновское рассеяние, которое играет ключевую роль в формировании спектров сильно замагниченных нейтронных звёзд (см., например, обзоры~\cite{Suleimanov:2007it,RCh09,Kuznetsov:2015}).
В частности, в работе~\cite{RCh09} 
выражение для сечения комптоновского рассеяния для случая, когда начальный и конечный электроны находятся на основном уровне Ландау, было получено с учётом дисперсии и перенормировки волновых функций фотонов. Одним из применений полученных результатов является решение уравнения переноса излучения в магнитосфере нейтронных звёзд. Подобная задача рассматривалась в работе~\cite{Mushtukov:2016}, где приводится уравнение переноса для дальнейшего точного решения.
В работе~\cite{Mushtukov:2022} эта задача была решена методом Монте-Карло и получены спектры для различных значений температуры и углов между направлением магнитного поля и импульсом фотона.

С другой стороны, есть уникальная возможность аналитически получить решение кинетического уравнения в окрестности резонанса для нахождения функции распределения фотонов двух возможных поляризаций в равновесной электрон-позитронной плазме в относительно сильном магнитном поле с учетом резонанса на виртуальном электроне.

\label{sec:preparation}
\section*{Решение кинетического уравнения вблизи резонанса}

Рассмотрим цилиндрическую колонку с осью, направленной вдоль магнитного поля (магнитное поле $B$ направлено вдоль оси $z$), содержащую равновесную $e^+e^-$ плазму при температуре $T=\text{const}\ll m$  и сильного магнитного поля $B\lesssim B_e$, через которую проходит стационарный поток фотонов. Будем предполагать, что поток фотонов описывается неравновесной функцией распределения. При этом основной вклад в концентрацию электронов и позитронов будет давать нулевой уровень Ландау. В таком случае  кинетическое уравнение будет иметь вид:
\begin{equation} 
	\begin{aligned}\label{Kineq:Gen}
	(\vec{n}, \vec{\bigtriangledown}_r f_\omega^{(\lambda)}(z, x))&=\sum_{\lambda'=1}^{2} \int dW_{\lambda\to\lambda'}\{f_{E'}(1-f_E)f^\lambda_{\omega'}(z, x)^{(\lambda')}(z, x) (1+f_\omega^{(\lambda)}(z, x))-
	\\
	&-f_E (1-f_{E'}) f_\omega^{(\lambda)}(1+f_{\omega'}(z, x)^{\lambda'})
%	\}\, .
	\end{aligned}
\end{equation}
Здесь $x = \cos \theta$, $x' = \cos \theta'$ ($\theta$, $\theta'$ -- угол между направлением магнитного поля и импульсами начального и конечного фотонов), $\lambda, \lambda'=1,2$ поляризационные состояния фотонов,  $\vec{n}$ -- единичный вектор, направленный вдоль импульса начального фотона, $f^{(\lambda)}_\omega(z, x)$ и $f^{(\lambda')}_{\omega'}(z, x')$ -- функции распределения начального и конечного фотонов с энергиями $\omega$ и $ \omega'$ соответственно, $dW_{\lambda \to \lambda'}$ -- коэффициент поглощения фотона, который может быть получен из результатов работ \cite{YarkovJETP:2017,YarkovJoPCS:2020}, $f_E = 1/[\exp(E/T)+1]$ и $f_{E'} = 1/[\exp(E'/T)+1]$ --  равновесные функции распределения начального и конечного электронов, $E, E'$~-- энергии начального и конечного электронов с нулевым химическим потенциалом, что характерно для магнитосфер нейтронных звёзд.
 

 Поскольку рассматривается нерелятивистская плазма с температурой $T\ll m$, то энергии начального и конечного фотона будут близкими друг к другу.
Раскладывая правую часть уравнения (\ref{Kineq:Gen}) по величине $\Delta\omega=\omega-\omega'\ll\omega$, где
	\begin{equation}
		\begin{aligned}
			\omega' =& \frac{1}{1-x'^2}\bigg(m+\omega(1-x x')-
			\\		
			-& \sqrt{(m+\omega(1-x x'))^2 - 2eB (1-x'^2)}\bigg)\simeq\omega_{\text{res}}\simeq eB/m\, ,
		\end{aligned}
	\end{equation}
мы можем воспользоваться методикой, развитой в работах \cite{Komp,Lubarsk}:
		\begin{equation}\label{eq:KinSerDeltaw}
					\begin{aligned}
			&\frac{\partial f^{(\lambda)}(z,x)}{\partial z} = \frac{1}{x} \sum_{\lambda'=1}^{2} \int_{-1}^{1}dx'\varphi_\omega^{\lambda\lambda'}(x,x') (f^{(\lambda')}(z,x')-f^{(\lambda)}(z,x)) +
			\\
			&- \frac{\Delta \omega}{T}\frac{1}{x} \sum_{\lambda'=1}^{2} \int_{-1}^{1}dx'\varphi_\omega^{\lambda\lambda'}(x,x') 
			\left[
			T \frac{\partial f_\omega^{(\lambda')}(z,x')}{\partial \omega} + f_\omega^{(\lambda')}(z,x')
			\right]+
			\\
			&+ \frac{1}{2}\frac{\Delta^2 \omega}{T^2}\frac{1}{x} \sum_{\lambda'=1}^{2} \int_{-1}^{1}dx'\varphi_\omega^{\lambda\lambda'}(x,x') 
			\bigg[
			T^2 \frac{\partial^2 f_\omega^{(\lambda')}(z,x')}{\partial \omega^2}+
			\\		
			&+ 2T \frac{\partial f_\omega^{(\lambda')}(z,x')}{\partial \omega^2}+f_\omega^{(\lambda')}(z,x')
			\bigg] 
			\,  .
			\end{aligned}
		\end{equation}

Рассматривая задачу вблизи резонанса $\omega=\omega_\text{res}\simeq eB/m$, функции $\varphi_\omega^{\lambda\lambda'}(x,x')$ могут быть получены из результата работы \cite{YarkovJoPCS:2020}
	\begin{equation}
		\begin{aligned}
		&\varphi_\omega^{11}(x,x')\simeq \rho \left(\frac{\beta^2}{m^2}\frac{1}{\Delta^-} +\frac{\left[\omega-eB/m\right]^2}{\Delta^+}\right)\, ,
		\\
		&\varphi_\omega^{12}(x,x')\simeq \rho \omega^2 \left(
		\frac{x^2}{\Delta^-} + \frac{\omega - eB/m}{2m^2 \Delta^+}
		\right)	\, ,
		\\
		&
		\varphi_\omega^{22}(x,x')\simeq \rho \frac{\omega^4}{4m^2}
		\left(
			 \frac{4m^4x^2 x'^2 }{(eB)^2 \Delta^-} + \frac{1}{\Delta^+}
		\right)\, ,
		\\
		&
				\varphi_\omega^{21}(x,x')\simeq \rho \omega^2
		\left(
		\frac{x^2}{\Delta^-} + \frac{\omega - eB/m}{2m \Delta^+}
		\right)\, ,
		\end{aligned}
	\end{equation}
где
\begin{equation}
\rho = \frac{n_e}{8\pi} e^4 \left(\frac{\sqrt{m^2+eB} + m}{\sqrt{m^2+eB}}\right)\simeq \frac{n_e}{2\pi}\, ,
\end{equation}
$n_e$ -- концентрация электронов,
	\begin{equation}
		\Delta^\pm = (\omega^2(1-x^2)+2 \omega m - 2\beta)^2 + \left(\frac{E_n{''} \Gamma^\pm}{2}\right)^2\simeq \left(\frac{E_n{''} \Gamma^\pm}{2}\right)^2 \, ,
	\end{equation}
	где $\Gamma^\pm$ -- полная ширина поглощения электрона для двух возможных поляризационных состояний \cite{KM_Book_2003}.

Решение уравнения (\ref{eq:KinSerDeltaw}) формально можно представить следующим образом 

\begin{equation}\label{FormSol}
\begin{aligned}
f_\omega^{(\lambda)}(z,x)&=f_{0\omega} 	e^{-{\cal \chi}_\omega^{(\lambda)}(x)\cdot z} + \frac{1}{x}\int_{0}^{z}dz'\int_{-1}^{1}dx' e^{-{\cal \chi}_\omega^{(\lambda)}(x)\cdot (z-z')}\times 
\\
&\times\varphi_\omega^{\lambda 	\lambda'}(x,x')F_\omega^{(\lambda')}(z',x')\, ,
\end{aligned}	
\end{equation}
где $f_{0\omega}=[\exp(\omega/T)-1]^{-1}$ -- равновесная функция распределения фотонов, а
\begin{equation}\label{eq:ThrF}
	\begin{aligned}
		&F_\omega^{(\lambda')}(z',x') = f_\omega^{(\lambda')}(z',x') - \frac{\Delta \omega}{T}
		\left[
		T \frac{\partial f_\omega^{(\lambda')}(z',x')}{\partial \omega} + f_\omega^{(\lambda')}(z',x')
		\right]+
		\\
		&+ \frac{1}{2}\frac{\Delta^2 \omega}{T^2}
		\left[
		T^2 \frac{\partial^2 f_\omega^{(\lambda')}(z',x')}{\partial \omega^2} + 2T \frac{\partial f_\omega^{(\lambda')}(z',x')}{\partial \omega^2}+f_\omega^{(\lambda')}(z',x')
		\right]\, ,
	\end{aligned}
\end{equation}
	\begin{equation}
{\cal \chi}_\omega^{(\lambda)}(x) \equiv \frac{1}{x} \int_{-1}^{1} dx' \left\{\varphi^{\lambda 1}_\omega(x,x')+\varphi^{\lambda 2}_\omega(x,x')\right\}\, .
\end{equation}

Далее раскладываем функцию распределения фотонов по полиномам Лежандра ${\cal P}_\ell(x)$ относительно переменной $x$:
\begin{equation}\label{eq:fSerLegender}
	f_\omega^{(\lambda)}(z, x) = \sum_{\ell = 0}^{\infty} A_\ell^{(\lambda)}(z,\omega) {\cal P}_\ell(x)\, ,
\end{equation}

Подставляя (\ref{eq:fSerLegender}) в (\ref{FormSol}) с учетом (\ref{eq:ThrF}) и применяя к полученному выражению преобразование Лапласа относительно переменной $z$, получим систему дифференциальных уравнений второго порядка относительно переменной $\omega$
\begin{equation}\label{eq:SysEqALap}
\begin{aligned}
\frac{2}{2 \ell +1} \overline{A}_\ell^{(\lambda)}(s,\omega) &= \int_{-1}^{1} \frac{f_{0\omega}^{(\lambda)}}{s+{\cal \chi}_\omega(x)}{\cal P}_\ell(x)dx+
\\
&+ \sum_{\ell'=0}^{\infty}\frac{1}{x}\int_{-1}^{1}dx'\int_{-1}^{1}dx \frac{P_\ell(x) P_{\ell'}(x')}{s+{\cal \chi}_\omega^{(\lambda)}(x)} \varphi_\omega^{\lambda\lambda'}(x,x') {\overline{\cal F}}_{\ell'}^{(\lambda')}(s,\omega)\, ,
\end{aligned}
\end{equation}
где 
\begin{equation}\label{eq:Flap}
\begin{aligned}
&{\overline{\cal F}}_{\ell'}^{(\lambda')}(s,\omega) = \overline{A}_{\ell'}^{(\lambda')}(s,\omega) - \frac{\Delta \omega}{T}
\left[
T \frac{\partial \overline{A}_{\ell'}^{(\lambda')}(s,\omega)}{\partial \omega} + \overline{A}_{\ell'}^{(\lambda')}(s,\omega)
\right]+
\\
&+ \frac{1}{2}\frac{\Delta^2 \omega}{T^2}
\left[
T^2 \frac{\partial^2 \overline{A}_{\ell'}^{(\lambda')}(s,\omega)}{\partial \omega^2} + 2T \frac{\partial \overline{A}_{\ell'}^{(\lambda')}(s,\omega)}{\partial \omega^2}+\overline{A}_{\ell'}^{(\lambda')}(s,\omega)
\right]\, ,
\end{aligned}
\end{equation}

\begin{equation}\label{lapA}
\overline{A}_{\ell'}^{(\lambda)}(s,\omega)=\int_{0}^{\infty} A^{(\lambda)}_{\ell'}(z,\omega) e^{-sz}dz\, ,
\end{equation}



Система уравнений (\ref{eq:SysEqALap}) с учетом (\ref{eq:Flap}) решает задачу о нахождении функции распределения фотонов. После применения обратного преобразования Лапласа, окончательная функция распределения фотонов поляризации $\lambda$ можно представить в следующем виде
\begin{equation}
f_\omega^{(\lambda)}(z, x) = \frac{1}{2\pi i}\sum_{\ell=0}^{\infty} 
P_\ell(x)\int_{\sigma-i\infty}^{\sigma + i \infty}ds \cdot e^{sz} \overline{A}^{(\lambda)}_\ell(s,x,\omega)\, ,
\end{equation}
где интеграл берется в комплексной плоскости по прямой $\text{Re} \, s = \sigma$
\section*{Заключение}
Рассмотрено решение кинетического уравнения для нахождения функции распределения фотонов двух возможных поляризаций в равновесной нерелятивистской $e^+e^-$ плазме в относительно сильном магнитном поле с учётом резонанса на виртуальном электроне.
Используя преобразование Лапласа и разложение функции распределения по полиномам Лежандра задача сведена к системе дифференциальных уравнений на коэффициенты этого разложения. Решение представлено в виде интеграла в комплексной плоскости.

Исследование выполнено при финансовой поддержке РФФИ в рамках научного проекта № 20-32-90068.
%============================= Fig. 1 ================================
%\begin{figure}[t]
%\begin{center}
%\includegraphics[width=127mm]{fig_1.eps}
%\includegraphics[width=127mm]{fig_1.eps}
%\vspace{-3mm}
%\caption{Figure caption}
%\end{center}
%\labelf{fig01}
%\vspace{-5mm}
%\end{figure}
%============================= Fig. 1 ================================

%\nocite{*}
%\bibliographystyle{pepan}
%\bibliography{maikbibl}
% Bibliography automatically generated via BibTeX (see template_bibtex.tex)

\begin{thebibliography}{10}
\def\selectlanguageifdefined#1{
\expandafter\ifx\csname date#1\endcsname\relax
\else\selectlanguage{#1}\fi}
\providecommand*{\href}[2]{{\small #2}}
\providecommand*{\url}[1]{{\small #1}}
\providecommand*{\BibUrl}[1]{\url{#1}}
\providecommand{\BibAnnote}[1]{}
\providecommand*{\BibEmph}[1]{\emph{#1}}
\ProvideTextCommandDefault{\cyrdash}{\hbox to.8em{--\hss--}}
\providecommand*{\BibDash}{\ifdim\lastskip>0pt\unskip\nobreak\hskip.2em\fi
\cyrdash\hskip.2em\ignorespaces}


\bibitem{D.Ryu:2012}
\selectlanguageifdefined{english}
\BibEmph{Ryu~D., Schleicher~D.~R.~G. Treumann~R.~A.Tsagas~C.~G., Widrow L.~M.} {Magnetic fields in the large-scale structure of the universe}~//
Space Science Reviews
\BibDash
\newblock 2012, \BibDash
\newblock V.~166, no.~1-4 \BibDash
\newblock P.~1-35.

\bibitem{Goldreich:1969}
\selectlanguageifdefined{english}
\BibEmph{Goldreich P., Julian W.~H.} {Pulsar electrodynamics}~//
Astrophys. J.
\BibDash
\newblock 1969, \BibDash
\newblock V.~157, \BibDash
\newblock P.~869-880.

\bibitem{Suleimanov:2007it}
\selectlanguageifdefined{english}
\BibEmph{Suleimanov V., Werner K.,} {Importance of Compton scattering for radiation spectra of isolated neutron stars with weak magnetic fields}~//
  Astron. Astrophys.
  \BibDash
\newblock 2007. \BibDash
\newblock P.~661--668.

\bibitem{RCh09}
\selectlanguageifdefined{english}
\BibEmph{Chistyakov, M. V., Rumyantsev, D. A.} {Compton effect in strongly magnetized plasm}~//
Int. J. Mod. Phys.
\BibDash
\newblock 2009. \BibDash
\newblock V.~24 \BibDash
\newblock P.~3995--4008.

\bibitem{Kuznetsov:2015}
\selectlanguageifdefined{english}
\BibEmph{Kuznetsov A. V., Rumyantsev D. A., Shlenev, D.
	M.} {Generalized two-point tree-level amplitude $j f \to 
	j'f'$ in a magnetized medium}~//
Int. J. Mod. Phys.
\BibDash
\newblock 2015. \BibDash
\newblock V.~A30, no.~11. \BibDash
\newblock P.~1550049.

\bibitem{Mushtukov:2016}
\selectlanguageifdefined{english}
\BibEmph{Mushtukov A. A., Nagirner D. I.,
	Poutanen J.} {Compton scattering S-matrix and cross section in strong
	magnetic field}~//
Phys. Rev.
\BibDash
\newblock 2016. \BibDash
\newblock V.~D93, no.~10. \BibDash
\newblock P.~105003.


\bibitem{Mushtukov:2022}
\selectlanguageifdefined{english}
\BibEmph{Mushtukov A.A., Markozov I.D., Suleimanov V.F., Nagirner D.I., Kaminke A.D., Potekhin, A.Y., Portegies Zwart S.} {Statistical features of multiple Compton scattering in a strong magnetic field}~//
Phys. Rev. D
\BibDash
\newblock 2022. \BibDash
\newblock V.~105, \BibDash
\newblock P.~103027.

\bibitem{YarkovJETP:2017}
\selectlanguageifdefined{english}
\BibEmph{D. A. Rumyantsev, D. M. Shlenev, A. A. Yarkov } {Resonances in Compton-Like scattering processes in an external magnetized medium}~//
JETP
\BibDash
\newblock 2017. \BibDash
\newblock V.~125, \BibDash
\newblock P.~410--419.

\bibitem{YarkovJoPCS:2020}
\selectlanguageifdefined{english}
\BibEmph{Chistyakov M.~V., Rumuyantsev D.~A., Yarkov A.~A. } {Effect of a strongly magnetized plasma on the resonant photon scattering process}~//
J. Phys.: Conf. Ser.
\BibDash
\newblock 2020. \BibDash
\newblock P.~1690 012015.

\bibitem{Komp}
\selectlanguageifdefined{english}
\BibEmph{Kompaneets A.S.} {The Establishment of Thermal Equilibium between Quanta and Electrons}~//
JETP
\BibDash
\newblock 1957. \BibDash
\newblock V.~4 no. 5, \BibDash
\newblock P.~876 -- 885.

\bibitem{Lubarsk}
\selectlanguageifdefined{russian}
\BibEmph{Любарский Ю. Э.} {Комптонизация в сверхсильном магнитном \\поле.~I}~//
Астрофизика
\BibDash
\newblock 1988. \BibDash
\newblock Том 28, \BibDash
\newblock Выпуск 1.

\bibitem{KM_Book_2003}
\selectlanguageifdefined{english}
\BibEmph{Kuznetsov A. V., Mikheev N. V.} {Electroweak processes in external electromagnetic fields}~//
Springer-Verlag
\BibDash
\newblock 2003, \BibDash
\newblock V.~120.


\end{thebibliography}
\end{document}
